% ---
% RESUMO
% ---
% resumo na língua vernácula (obrigatório)
\begin{resumo} %% AQUI COMEÇA A PÁGINA DE RESUMO

A atribuição de disciplinas a professores em instituições de ensino superior é um problema combinatório recorrente, classificado como NP-difícil por sua relação com o \textit{Generalized Assignment Problem}, cuja complexidade cresce com o número de docentes, disciplinas e restrições operacionais envolvidas, tornando métodos exatos baseados em Programação Linear Inteira (PLI) potencialmente inviáveis em tempo computacional para instâncias de maior porte, ao passo que abordagens puramente heurísticas não oferecem garantias de qualidade. Este trabalho propõe, implementa e avalia uma matheurística baseada em Large Neighborhood Search (LNS) com reotimização por sub-MIP, em um esquema do tipo \textit{fix-and-optimize}, integrada como heurística primal ao solver SCIP 7.0, para o Problema da Distribuição de Disciplinas (PDD), modelado como um programa linear inteiro binário cuja função objetivo maximiza a qualidade da atribuição com base em preferências docentes e compatibilidade de área, sujeito a restrições de cobertura obrigatória de todas as disciplinas, carga horária mínima anual e máxima semestral por professor. A heurística LNS opera como mecanismo de intensificação integrado ao Branch-and-Bound (B\&B): a partir de uma solução incumbente, identifica as atribuições ativas, ordena os candidatos segundo uma métrica de versatilidade docente e libera uma fração controlada das variáveis para reotimização em um subproblema auxiliar (sub-MIP) com limite de tempo, incorporando a solução resultante ao solver principal caso represente melhoria. Três parâmetros internos --- taxa de destruição, sentido de ordenação dos candidatos e tempo máximo do sub-MIP --- foram calibrados por meio de uma abordagem incremental \textit{one-factor-at-a-time}, utilizando subconjuntos de 10 instâncias extraídos de um conjunto de 42 instâncias sintéticas, com limite de tempo global de 400 segundos por instância. A combinação selecionada foi de 75\% de destruição, ordenação crescente e 200 segundos de sub-MIP. No comparativo final entre B\&B puro, LNS, heurística construtiva \textit{Greedy Randomized Adaptive Search Procedure} (GRASP) e variante híbrida GRASP|LNS, foram avaliadas todas as 42 instâncias sintéticas e uma instância real. Nenhuma abordagem dominou todas as métricas: o B\&B puro obteve o menor tempo total e o menor gap médio; GRASP|LNS apresentou o menor total de nós restantes; e B\&B puro e LNS empataram com 32 das 43 instâncias resolvidas até a otimalidade. Na instância real, LNS e B\&B puro provaram a otimalidade, com pequena vantagem temporal para o LNS. Os resultados mostram que a contribuição das heurísticas depende da métrica e da instância, destacando a redução da árvore proporcionada pela variante híbrida e o desempenho competitivo do LNS no caso real, sem evidência de superioridade global sobre o método exato.

\noindent
    {\sloppy \textbf{Palavras-chave}: Distribuição de Disciplinas. Programação Linear Inteira. Branch-and-Bound. Matheurística. Large Neighborhood Search. Fix-and-Optimize. SCIP.\par}\fussy
\end{resumo} %AQUI TERMINA A PÁGINA DE RESUMO
